\chapter{核心方法与实验分析}
\label{chap:method}

\section{问题形式化描述}
我们将论文协同修订建模为带显式人工确认门控的状态迁移序列，目标损失函数如公式~\eqref{eq:loss}~所示：
\begin{equation}
  \mathcal{L}(\theta) = \sum_{i=1}^{N} \ell(f_\theta(x_i), y_i) + \lambda \|\theta\|_2^2
  \label{eq:loss}
\end{equation}

\section{模型优化过程}
参数迭代更新与收敛判定由多行独立编号公式~\eqref{eq:update}~与公式~\eqref{eq:stop}~给出：
\begin{align}
  \theta^{(t+1)} &= \theta^{(t)} - \eta \nabla_\theta \mathcal{L}(\theta^{(t)}) \label{eq:update} \\
  \Delta_t &= \|\theta^{(t+1)} - \theta^{(t)}\|_2 < \varepsilon \label{eq:stop}
\end{align}

\section{实验设计与结果分析}
\subsection{实验数据集与对比结果}
表~\ref{tab:exp}~展示了不同方法在中文（CH）与英文（EN）基准数据集上的对比实验结果（标准多级表头三线表示例）。

\begin{table}[htbp]
  \centering
  \caption{不同基线方法在基准数据集上的性能对比}
  \label{tab:exp}
  \begin{tabular}{lcccc}
    \toprule
    \multirow{2}{*}{\textbf{方法}} & \multicolumn{2}{c}{\textbf{CH 数据集}} & \multicolumn{2}{c}{\textbf{EN 数据集}} \\
    \cmidrule(lr){2-3} \cmidrule(lr){4-5}
    & Acc & $F_1$ & Acc & $F_1$ \\
    \midrule
    基线方法 A\cite{he2016deep} & 0.869 & 0.868 & 0.845 & 0.798 \\
    基线方法 B\cite{devlin2019bert} & 0.907 & 0.913 & 0.881 & 0.820 \\
    \textbf{本文方法（Ours）} & \textbf{0.951} & \textbf{0.942} & \textbf{0.917} & \textbf{0.910} \\
    \bottomrule
  \end{tabular}
\end{table}

\section{本章小结}
本章提出了核心优化框架，并通过对比实验验证了方法的有效性与稳健性。
